\documentclass{article}
\usepackage{graphicx}
\usepackage{amsmath, amssymb, amsthm}
\usepackage[T1]{fontenc}
\usepackage[utf8]{inputenc}
\usepackage{hyperref}
\usepackage{enumitem}
\usepackage{geometry}
\geometry{margin=1in}

\title{tagd Math}
\author{Isaac N Colley}
\date{April 2026}

\newtheorem{definition}{Definition}[section]
\newtheorem{proposition}{Proposition}[section]
\newtheorem{remark}{Remark}[section]

\begin{document}

\maketitle

\section{Introduction}

This document describes the mathematical structures already present in the
\texttt{tagd} semantic-relational system and the TAGL language specification.
Its purpose is not to invent a new mathematical model, but to state, as
precisely as possible, the mathematics implicit in the current source and
specification.

The central claim is that a tagspace is not merely a collection of labels.
It has mathematical structure. In particular, a tagspace exhibits at least
the following:

\begin{itemize}
    \item a set of tags
    \item a rooted tree induced by subordinate relations
    \item a rank-based containership relation
    \item a partial order induced by rank-prefix containment
    \item additional labeled relations on nodes through predicates
\end{itemize}

These structures support the broader thesis that in \texttt{tagd},
information is structural: meaning is determined not only by nominal labels,
but by position in the hierarchy and by relations attached to that position.

\section{Preliminaries}

\subsection{Tags as a Set}

\begin{definition}
Let $T$ be the set of all tags in a tagspace.
\end{definition}

Each tag has two forms of identity:

\begin{enumerate}[label=\arabic*.]
    \item \textbf{Nominal identity}: a globally unique UTF-8 text label.
    \item \textbf{Structural identity}: its subordinate relation and its
    predicate list.
\end{enumerate}

Example TAGL statements:

\begin{verbatim}
>> mammal _type_of animal

>> dog _is_a mammal
\end{verbatim}

Here, \texttt{dog}, \texttt{mammal}, and \texttt{animal} are all elements of
$T$.

\subsection{Subordinate Relations}

\begin{definition}
A subordinate relation is a binary relation on tags:
\[
\mathrel{\_sub} \subseteq T \times T
\]
where $(a,b) \in \mathrel{\_sub}$ means that tag $a$ is subordinate to tag $b$.
\end{definition}

In TAGL surface syntax, subordinate relations may be written using
particular sub-relators such as \texttt{\_sub}, \texttt{\_is\_a}, or
\texttt{\_type\_of}. The general structural role is the same: a subject is
placed under a super-object in the tagspace hierarchy.

Example:

\begin{verbatim}
>> dog _is_a mammal
\end{verbatim}

\section{Tagspace as a Rooted Tree}

\subsection{Root}

\begin{definition}
There exists a distinguished root tag
\[
\_entity \in T
\]
which is axiomatic.
\end{definition}

The root relation is:

\begin{verbatim}
_entity _sub _entity
\end{verbatim}

This makes \texttt{\_entity} the unique root of the tagspace.

\subsection{Parent Function}

Because each tag has exactly one subordinate relation, each non-root tag has
exactly one parent in the hierarchy.

\begin{definition}
Define the parent function
\[
\mathrm{parent}: T \to T
\]
such that
\[
\mathrm{parent}(a) = b
\]
whenever $a$ is subordinate to $b$.
\end{definition}

\begin{proposition}
The subordinate structure of a tagspace forms a rooted tree.
\end{proposition}

\begin{proof}
Each tag has exactly one subordinate relation, and all tags ultimately
descend from the distinguished root $\_entity$. Therefore the hierarchy has
a unique root and unique parent assignment for each node, which is exactly
the defining shape of a rooted tree.
\end{proof}

\subsection{Example}

\begin{verbatim}
>> physical_object -^ _entity
>> living_thing _type_of physical_object
>> animal _type_of living_thing
>> mammal _type_of animal
>> dog _is_a mammal
\end{verbatim}

This yields the tree fragment

\[
\_entity \leftarrow physical\_object \leftarrow living\_thing \leftarrow
animal \leftarrow mammal \leftarrow dog
\]

where arrows point from child to parent.

\section{Rank and Containership}

\subsection{Rank Function}

In \texttt{tagd}, each tag has a rank. A rank is a UTF-8 byte string whose
prefixes encode containership in the tree.

\begin{definition}
Define a rank function
\[
\mathrm{rank}: T \to \Sigma^{*}
\]
where $\Sigma^{*}$ is the set of finite UTF-8 strings.
\end{definition}

The intended semantics are:

\begin{itemize}
    \item the rank of a parent is a prefix of the rank of each descendant
    \item rank substrings define subtrees
    \item lexicographic collation of rank strings preserves hierarchy
\end{itemize}

\subsection{Containership Relation}

\begin{definition}
Define containership on tags by
\[
b \preceq a \quad \Longleftrightarrow \quad \mathrm{rank}(b)
\text{ is a prefix of } \mathrm{rank}(a)
\]
\end{definition}

This means that $b$ contains $a$ in the hierarchical sense: $a$ lies in the
subtree rooted at $b$.

\begin{remark}
This is the mathematical form of the source-level intuition:
a parent rank prefixes all ranks in its subtree.
\end{remark}

\subsection{Example}

Suppose the ranks are:

\[
\mathrm{rank}(\_entity)=\epsilon,\quad
\mathrm{rank}(mammal)=1.2.4,\quad
\mathrm{rank}(dog)=1.2.4.7
\]

Then $\mathrm{rank}(mammal)$ is a prefix of $\mathrm{rank}(dog)$, so
$mammal \preceq dog$.

\section{Tagspace as a Partially Ordered Set}

\subsection{Specialization Order}

There are two equivalent ways to orient the order. For language design,
the most useful one is usually the specialization order.

\begin{definition}
Define
\[
a \leq b \quad \Longleftrightarrow \quad \mathrm{rank}(b)
\text{ is a prefix of } \mathrm{rank}(a)
\]
\end{definition}

Interpretation:

\begin{itemize}
    \item $a \leq b$ means $a$ is subordinate to $b$
    \item $a$ is more specific than $b$
    \item $b$ is more general than $a$
\end{itemize}

\begin{proposition}
$(T,\leq)$ is a partially ordered set.
\end{proposition}

\begin{proof}
\textbf{Reflexive:} Every rank is a prefix of itself, so $a \leq a$.

\textbf{Transitive:} If $\mathrm{rank}(b)$ is a prefix of
$\mathrm{rank}(a)$ and $\mathrm{rank}(c)$ is a prefix of
$\mathrm{rank}(b)$, then $\mathrm{rank}(c)$ is a prefix of
$\mathrm{rank}(a)$, so $a \leq c$.

\textbf{Antisymmetric:} If $a \leq b$ and $b \leq a$, then each rank is a
prefix of the other, hence the ranks are equal. In a tagspace, equal
structural position implies the same node in the tree, so $a=b$.
\end{proof}

\subsection{Interpretation}

This partial order is not added from outside. It is already present in the
system through rank-prefix containment. It formalizes the intuition that
the tagspace expresses a hierarchy of specialization and generalization.

\section{Subtrees and Set-Theoretic Structure}

\subsection{Subtrees as Sets}

Given any tag $x \in T$, define its subtree by

\[
\mathrm{Subtree}(x) = \{ y \in T \mid x \preceq y \}
\]

This is the set of all descendants of $x$, including $x$ itself.

\subsection{Examples}

If the tagspace contains:

\begin{verbatim}
>> animal _type_of living_thing
>> mammal _type_of animal
>> bird _type_of animal
>> dog _is_a mammal
>> cat _is_a mammal
\end{verbatim}

then

\[
\mathrm{Subtree}(animal)
=
\{animal, mammal, bird, dog, cat, \dots \}
\]

and

\[
\mathrm{Subtree}(mammal)
=
\{mammal, dog, cat, \dots \}
\]

\subsection{Intersection Behavior}

For two tags $x,y \in T$, the intersection
\[
\mathrm{Subtree}(x) \cap \mathrm{Subtree}(y)
\]
is either:

\begin{itemize}
    \item empty, if the subtrees are disjoint
    \item $\mathrm{Subtree}(x)$, if $x$ lies below $y$
    \item $\mathrm{Subtree}(y)$, if $y$ lies below $x$
\end{itemize}

This follows from the fact that a rooted tree has no cycles and each node
has a unique path to the root.

\section{A Topological View of Tagspaces}

\subsection{Motivation}

The tagspace is already a set, a rooted tree, and a partially ordered set.
This suggests a further mathematical question:

\begin{quote}
Can the containership structure of a tagspace be described topologically?
\end{quote}

The answer is yes. A natural topology arises from the order structure.

\subsection{Alexandrov Topology}

Any partial order induces an Alexandrov topology.

\begin{definition}
Let $(T,\leq)$ be the specialization order above. A subset
$U \subseteq T$ is open if
\[
x \in U \text{ and } x \leq y \implies y \in U
\]
for all $x,y \in T$.
\end{definition}

These are the upward-closed sets with respect to the order $\leq$.
Because $\leq$ means ``more specific than,'' an upward-closed set contains
everything more general than any tag it already contains.

Equivalently, if one prefers to work with the containership order
$\preceq$, one can use the dual presentation in which descendant-closed
subtrees play the role of basic open sets.

\subsection{Why This Fits Tagd}

The topology is not arbitrary. It is induced by the same structural facts
already present in the source:

\begin{itemize}
    \item ranks define a prefix order
    \item prefix order defines containment
    \item containment defines subtrees
    \item subtree behavior is closed under the natural set operations of the tree
\end{itemize}

\subsection{Topological Intuition}

In a tagspace, generalization and specialization propagate structurally.

Example:

\begin{verbatim}
>> animal _type_of living_thing
>> mammal _type_of animal
>> dog _is_a mammal
\end{verbatim}

Then $dog \leq mammal \leq animal \leq living\_thing$.

This means that moving upward in the order moves toward more general
structure. In this sense, the topology encodes the taxonomic shape of the
tagspace.

\section{Hard Tags as a Universal Invariant Substructure}

\subsection{Hard Tags}

Let
\[
H \subseteq T
\]
be the set of hard tags.

These are the tags defined axiomatically by the system, including tags
such as:

\begin{itemize}
    \item \texttt{\_entity}
    \item \texttt{\_sub}
    \item \texttt{\_rel}
    \item \texttt{\_number}
    \item \texttt{\_interrogator}
\end{itemize}

\subsection{Shared Base}

Every tagspace contains the same hard tag base. Therefore $H$ is a common
invariant substructure across all tagspaces.

\begin{remark}
The hard tag base is the structural common ground that makes all tagspaces
mutually legible. Any extension of a tagspace is built on top of the same
axiomatic base.
\end{remark}

\subsection{Example}

A bootstrap dump shows a fixed base such as:

\begin{verbatim}
_entity
_sub
_is_a
_type_of
_rel
_has
_can
_number
_integer
_float
_interrogator
_what
_search
\end{verbatim}

No matter how different two tagspaces become, this hard-tag substructure
remains shared.

\section{Predicate Relations as Additional Structure}

\subsection{Predicates}

A predicate relation is not part of the vertical containment tree.
Instead, it adds labeled relational structure to nodes.

\begin{definition}
A predicate is a relation of the form
\[
(r, O)
\]
where $r$ is a relator and $O$ is an object list, applied to a subject.
\end{definition}

A predicate with modifier may be represented more explicitly as
\[
(r, o, m)
\]
where $r$ is the relator, $o$ the object, and $m$ the modifier.

\subsection{Example}

\begin{verbatim}
>> dog _has legs = 4
>> dog _can bark
\end{verbatim}

These do not change the position of \texttt{dog} in the tree.
They add horizontal relational structure to the node \texttt{dog}.

\subsection{Mathematical View}

A tagspace therefore has at least two interacting layers:

\begin{enumerate}[label=\arabic*.]
    \item a vertical tree / order / topology given by subordinate relations
    \item a horizontal labeled relational layer given by predicates
\end{enumerate}

This is why the system is properly called semantic-relational. The meaning
of a tag is not exhausted by its nominal name; it is expressed by both
hierarchical placement and attached relations.

\section{Structural Identity}

\subsection{Two Forms of Identity}

A tag has:

\begin{enumerate}[label=\arabic*.]
    \item nominal identity
    \item structural identity
\end{enumerate}

Nominal identity is its globally unique UTF-8 label.

Structural identity is determined by:

\begin{itemize}
    \item its subordinate relation
    \item its rank
    \item its predicate list
\end{itemize}

\subsection{Interpretation}

This is a structuralist view of information. A tag is not meaningful in
isolation. Its meaning is determined by where it sits in the tagspace and
how it relates to other tags.

A concise formula is:

\[
\mathrm{meaning}(x) =
\mathrm{position}(x) + \mathrm{relations}(x)
\]

where ``position'' is given by subordinate structure and rank, and
``relations'' are given by predicates.

\section{Canonical Representation}

\subsection{Text as Structural Serialization}

A tagspace has a canonical textual representation through \texttt{dump()}.

This means that the full structure of the tagspace is representable in
deterministic TAGL text.

\subsection{Interpretation}

A tagspace is not merely stored as text. Rather, the text is a canonical
serialization of the underlying mathematical structure.

This is important because it allows:

\begin{itemize}
    \item deterministic inspection
    \item canonical comparison
    \item hashable structural snapshots
    \item faithful reload of the same structure
\end{itemize}

\section{Conclusion}

The \texttt{tagd} system already contains a rich mathematical foundation.

A tagspace may be described as:

\begin{itemize}
    \item a set of tags
    \item a rooted tree
    \item a rank-based containership hierarchy
    \item a partially ordered set under specialization
    \item an induced Alexandrov topology
    \item an invariant hard-tag substructure
    \item a relational layer of predicates attached to nodes
\end{itemize}

These are not external metaphors. They are mathematical descriptions of
structure already present in the system.

This supports the broader thesis that in \texttt{tagd} and TAGL,
information is structural.

\section{Future Work}

Natural continuations of this document include:

\begin{itemize}
    \item formal treatment of query semantics as set and order operations
    \item formal treatment of predicate algebra
    \item morphisms between tagspaces
    \item merge operations as structure-preserving transformations
    \item closure, inheritance, and specialization in the topological model
\end{itemize}

\end{document}

